\scanpage{440}
\mspage{7}
% archive sheet number: 439

Le théorème suivant doit être vrai, et résulter de \unsure{mes} fourbis généraux.

\uline{Théorème} Si le \deleted{\unsure{morphisme}} carré $(*)$ est \uline{cartésien} \inserted{\unsure{de|dans} $(\mathrm{Cat})$},

\noindent\begin{tabular}{@{}|p{.93\linewidth}@{}}
$\unsure{\gamma|r}$ \deleted{l'identité} un iso, et si on suppose que $f$ soit 1-lisse,
ou $p$ 1-propre, alors \unsure{le|la} \deleted{\unsure{signe}} $T'$-foncteur de
changement de base est une \unsure{équivalence} de \inserted{$T'$-}cat..
\end{tabular}

Les ``fourbis'' généraux seraient \deleted{\unsure{serait}} en termes \deleted{\unsure{termes}} d'un dérivateur $\unsure{\mathbb D|\mathrm D}$ ad hoc, celui des 1-types d'homotopie relatifs.. J'y reviendrai par la suite.

On se pose la question quand le 2-foncteur associé à \deleted{\unsure{$f$}} $f\colon S\longrightarrow T$

\[
f^*\colon\uuline{\mathrm{Grp}}(T)\unsure{\mathrel{-}|\longrightarrow}\uuline{\mathrm{Grp}}(S)
\]

est 0-fidèle, 1-fidèle, 2-fidèle, \deleted{\unsure{i.e.}} \deleted{Cela} Si $G,H$ sont \deleted{des} \unsure{gr|g} fibrés en groupoïdes sur $T$, on considère

\margin{La \unsure{$i$-fidélité|$i$-infidélité} de $f^*$ ($i\in\{\deleted{\unsure{n}},0,1,2\}$) équivaut à celle des foncteurs (obtenus \deleted{\unsure{à|en}} en prenant $G=e_T$) $\uline{\unsure{\Gamma|\mathscr T}}(H/T)\longrightarrow\uline{\Gamma}(f^*H/S)$.}
% slanted

\[
\uline{\mathrm{Hom}}_{\uuline{\mathrm{Grp}}(T)}(G,H)\unsure{\mathrel{-}|\longrightarrow}
\uline{\mathrm{Hom}}_{\uuline{\mathrm{Grp}}(S)}(f^*G,f^*H),
\]

et suivant que ce foncteur est toujours fidèle, ou pleinement fidèle, ou une équivalence de cat on dit que $f^*$ est 0-, 1-, ou 2-fidèle. Dans le cas 2-fidèle, cela \deleted{signifie} \inserted{implique} que si $K$ est un fibré en groupoïdes sur $S$, s'il est équivalent à un $f^*G$, alors $G$ est déterminé à \inserted{$T$-}équiv. près, elle \unsure{$=$|même} déterminée à iso unique près, \deleted{Si ça marche pour \uline{tt} $K$, je dirai que $f^*$ est 3-fidèle. On veut des critères \unsure{commodes} \deleted{pour} \deleted{\unsure{que ce soit}} de la $i$-fidélité.}
